\title{Group Sequence Policy Optimization}

\begin{abstract}
We present Group Sequence Policy Optimization (GSPO), a robust, efficient, and high-performing reinforcement learning approach designed for training large language models. In contrast to earlier methods that utilize token-level importance ratios, GSPO operates by establishing the importance ratio on entire sequence probabilities and applies sequence-level procedures for clipping, reward assignment, and optimization. Our empirical results indicate that GSPO provides enhanced training efficiency and achieves better performance relative to the GRPO method, exhibits greater stability in Mixture-of-Experts (MoE) RL scenarios, and can streamline RL infrastructure development. These advantages of GSPO have played a significant role in advancing the performance of the most recent Qwen3 models.
\end{abstract}

\section{Introduction}
\label{sec:intro}

Reinforcement learning (RL) has become a central approach for advancing the capabilities of language models at scale \citep{o1,dpsk-r1,qwq32b,qwen3}. By employing large-scale RL, these models acquire the ability to address complex challenges—including high-level mathematics and programming—by engaging in more intricate and extended forms of reasoning.

Achieving further scaling of RL via increased computational resources critically depends on ensuring stable and resilient training processes. Nevertheless, leading RL algorithms, such as GRPO \citep{grpo}, are plagued by significant instability when deployed on extremely large language models; this frequently leads to drastic and irreparable failures, known as model collapse \citep{qwen3, minimax-m1}. Such instability poses a major barrier to enhancing language model performance through ongoing RL training.

In this study, we demonstrate that the instability observed in GRPO arises from a fundamental misuse and subsequent invalidity of importance sampling weights within its framework. This flaw generates training noise with high variance, which accumulates as response sequences grow longer and is exacerbated by the use of a clipping mechanism, leading ultimately to the collapse of the model.

To overcome these intrinsic shortcomings, we introduce \textbf{Group Sequence Policy Optimization (GSPO)}, a novel reinforcement learning method tailored for training large language models. GSPO's primary contribution is its rigorous and theoretically sound formulation of the importance ratio, leveraging sequence likelihood in accordance with foundational importance sampling principles \citep{click}. Furthermore, GSPO determines normalized rewards as the advantage values across multiple responses per query, thereby maintaining consistency between sequence-level reward assignment and the optimization process.

Our experiments provide clear evidence that GSPO significantly outperforms GRPO in terms of training stability, efficiency, and overall effectiveness. Importantly, GSPO naturally addresses stability issues that commonly occur when applying RL to large Mixture-of-Experts (MoE) architectures, thereby obviating the requirement for intricate stabilization mechanisms. This simplification presents opportunities for streamlining RL infrastructure. As a result, the adoption of GSPO has led to marked enhancements in the performance of the most recent Qwen3 models. We anticipate that GSPO will serve as a dependable and scalable foundation for further progress in large-scale RL applications with language models.

\section{Preliminaries}

\paragraph{Notation}  
Throughout this work, we consider an autoregressive language model, whose parameters are represented by $\theta$, as a policy denoted by $\pi_\theta$. The variable $x$ stands for a query, and we represent the collection of all queries with $\mathcal{D}$. For any query $x$ and its corresponding response $y$, the conditional probability assigned by the policy $\pi_\theta$ is written as $\pi_\theta (y | x)=\prod_{t=1}^{|y|} \pi_\theta (y_t | x, y_{<t} )$, where $|y|$ indicates the length of $y$ in tokens. Each query-response pair $(x, y)$ may be evaluated by a verifier $r$, producing a reward value $r(x, y) \in [0, 1]$.

\paragraph{Proximal Policy Optimization (PPO)} Using samples generated from the old policy $\pi_{\theta_\text{old}}$, PPO \citep{ppo} constrains the policy update within a proximal region of the old policy through the clipping mechanism. Specifically, PPO employs the following objective for policy optimization (we omit the KL regularization term hereinafter for brevity, as it is not the focus of this paper):

\begin{align}
\mathcal{J}_\text{PPO}(\theta) = \mathbb{E}_{ x \sim \mathcal{D},\, y \sim \pi_{\theta_\text{old}}( \cdot | x) }
\left[ \frac{1}{|y|} \sum_{t=1}^{|y|} 
\min \left( w_{t}(\theta) \widehat{A}_{t},  \, \mathrm{clip} \left( w_{t}(\theta), 1 - {\varepsilon}, 1 + {\varepsilon}\right) \widehat{A}_{t} \right)
\right],
\end{align}

Here, the token importance ratio $w_{t}(\theta)$ is given by $
w_{t}(\theta) = \frac{ \pi_{\theta} (y_{t} | x, y_{<t}) }{ \pi_{\theta_\text{old}} (y_{t} | x,y_{<t})} $, the advantage term $\widehat{A}_{t}$ corresponding to $y_t$ is inferred by a separate value model, and $\varepsilon$ serves as the clipping parameter for the importance ratios.

A principal limitation of PPO in practical implementations is its strong dependence on the accuracy of the value model. Often, the value model is constructed with an architecture similar in scale to that of the policy model, which imposes significant computational and memory overhead. The algorithm’s performance further depends on the fidelity of its value predictions. Not only is learning a dependable value model intrinsically difficult, but scaling its accuracy to accommodate extended outputs and more intricate task demands becomes increasingly problematic.

\paragraph{Group Relative Policy Optimization (GRPO)} GRPO \citep{grpo} bypasses the need for the value model by computing the relative advantage of each response within a group of responses to the same query. Specifically, GRPO optimizes the following objective:

\begin{align}
\label{equ:grpo}
\mathcal{J}_\text{GRPO}(\theta) = \mathbb{E}_{ x \sim \mathcal{D},\, \{y_i\}_{i=1}^G \sim \pi_{\theta_\text{old}}( \cdot | x) }
\left[ \frac{1}{G} \sum_{i=1}^{G} \frac{1}{|y_i|} \sum_{t=1}^{|y_i|} 
\min \left( w_{i,t}(\theta) \widehat{A}_{i,t},  \, \mathrm{clip} \left( w_{i,t}(\theta), 1 - {\varepsilon}, 1 + {\varepsilon}\right) \widehat{A}_{i,t} \right)
\right],
\end{align}

where $G$ is the number of generated responses to each query $x$ (i.e., the group size), and the importance ratio $w_{i,t}(\theta)$ and advantage $\widehat{A}_{i,t}$ of token $y_{i,t}$ are:

\begin{align}
    w_{i,t}(\theta)=\frac{ \pi_{\theta} (y_{i,t} | x, y_{i,<t}) }{ \pi_{\theta_\text{old}} (y_{i,t} | x,y_{i,<t})},\quad\ 
    \widehat{A}_{i,t} = \widehat{A}_{i} = \frac{r(x, y_i) - \mathrm{mean} \left( \{ r(x, y_i) \}_{i=1}^G \right) }{ \mathrm{std} \left( \{ r(x, y_i) \}_{i=1}^G \right) },
\end{align}

respectively, where all the tokens in $y_i$ share the same advantage as $\widehat{A}_{i}$.

\section{Motivation}
\label{sec:motivation}

As models increase in size, adopt more sparse architectures such as Mixture-of-Experts, and generate longer responses, employing larger rollout batch sizes becomes crucial for efficient hardware usage during reinforcement learning. To enhance sample efficiency, it is common to subdivide these large rollout batches into several mini-batches for the purpose of performing gradient updates. However, this approach inherently results in an off-policy learning scenario, where samples $y$ are drawn from a previous policy $\pi_{\theta_\text{old}}$ rather than from the current policy $\pi_\theta$ undergoing training. This condition underlies the requirement for the clipping mechanisms in both PPO and GRPO, which serve to filter out samples that are excessively "off-policy" during the gradient estimation process.

While mechanisms like clipping aim to manage this off-policy discrepancy, we identify a more fundamental issue in GRPO: \textit{its objective is ill-posed}. This problem becomes particularly acute when training large models on long-response tasks, leading to catastrophic model collapse. The ill-posed nature of the GRPO objective stems from a misapplication of importance sampling weights. The principle of importance sampling is to estimate the expectation of a function $f$ under a target distribution $\pi_\text{tar}$ by re-weighting samples drawn from a behavior distribution $\pi_\text{beh}$:

\begin{align}
\mathbb{E}_{z \sim \pi_\text{tar}} \left[ f(z) \right] = \mathbb{E}_{z \sim \pi_\text{beh}} \left[ \frac{ \pi_\text{tar} (z) }{ \pi_\text{beh} (z)} \, f(z) \right].
\end{align}

Importantly, this approach depends on aggregating results from a large number of samples ($N \gg 1$) drawn from the behavioral policy $\pi_\text{beh}$, so that the importance weighting $\frac{ \pi_\text{tar} (z) }{ \pi_\text{beh} (z)}$ can accurately rectify any divergence between distributions.

Conversely, GRPO incorporates the importance weight $\frac{ \pi_{\theta} (y_{i,t} | x, y_{i,<t}) }{ \pi_{\theta_\text{old}} (y_{i,t} | x,y_{i,<t})}$ at every token timestep $t$. This calculation is predicated upon only a single observation $y_{i,t}$ from each next-token distribution $\pi_{\theta_\text{old}}( \cdot | x, y_{i, <t})$, which undermines its ability to correct the distribution shift as intended. Rather than serving its corrective purpose, this introduces substantial variance into the gradient estimates, with such noise accumulating over extended output sequences and further amplified by the effects of the clipping procedure. Our empirical results indicate that this scenario can culminate in irreversible model collapse. After such collapse, attempts to recover—even by rolling back to prior checkpoints and meticulously adjusting hyperparameters (including clipping ranges), increasing output lengths, or changing the RL queries—are unsuccessful.

This finding highlights a deep-seated limitation within the design of GRPO. The breakdown of token-level importance weighting reveals a fundamental tenet: \textbf{the granularity of the optimization objective must correspond to that of the reward function}. Given that the reward is allocated at the entire sequence level, off-policy corrections at the token level appear inherently inadequate. These considerations motivate abandoning the token-wise approach in favor of leveraging importance weighting and conducting optimization at the \textit{sequence level}.

\section{Algorithm}

\subsection{GSPO: Group Sequence Policy Optimization}

Although the token-level importance weight $\frac{ \pi_{\theta} (y_{i,t} | x, y_{i,<t}) }{ \pi_{\theta_\text{old}} (y_{i,t} | x,y_{i,<t})}$ presents challenges within GRPO, our analysis reveals that, for language generation tasks, the \textit{sequence-level} importance weight $\frac{ \pi_{\theta} (y | x) }{ \pi_{\theta_\text{old}} (y | x)}$ is theoretically well-founded. Specifically, this ratio quantifies the extent to which a response $y$ generated from $\pi_{\theta_\text{old}} (\cdot | x)$ diverges from the distribution defined by $\pi_{\theta} (\cdot | x)$. Consequently, it is well-suited for reflecting sequence-level rewards and provides a meaningful metric for the operation of the clipping mechanism.

Based on this straightforward observation, we propose the \textbf{Group Sequence Policy Optimization (GSPO)} algorithm. GSPO employs the following sequence-level optimization objective:

\begin{align}
\mathcal{J}_\text{GSPO} (\theta) =
\mathbb{E}_{ x \sim \mathcal{D},\, \{y_i\}_{i=1}^G \sim \pi_{\theta_\text{old}}( \cdot | x) }
\left[ 
\frac{1}{G} \sum_{i=1}^{G}
\min \left( s_{i}(\theta)  \widehat{A}_{i},  \, \mathrm{clip} \left( s_{i}(\theta), 1 - {\varepsilon}, 1 + {\varepsilon} \right) \widehat{A}_{i} \right) 
\right],
\label{equ:gspo}
\end{align}

where we adopt the group-based advantage estimation:

\begin{align}
\widehat{A}_{i} = \frac{r(x, y_i) - \mathrm{mean} \left( \{ r(x, y_i) \}_{i=1}^G \right) }{ \mathrm{std} \left( \{ r(x, y_i) \}_{i=1}^G \right) },
\end{align}

and define the importance ratio $s_{i}(\theta)$ based on sequence likelihood \citep{click}:

\begin{align}
s_{i}(\theta) = \left( \frac{ \pi_{\theta} (y_i | x) }{ \pi_{\theta_\text{old}} (y_i | x)} \right)^{\frac{1}{|y_i|}}
=
\exp \left( \frac{1}{|y_i|} \sum_{t=1}^{|y_i|} \log \frac{ \pi_{\theta} (y_{i,t} | x, y_{i,<t}) }{ \pi_{\theta_\text{old}} (y_{i,t} | x,y_{i,<t})} \right).
\end{align}

Consequently, GSPO implements clipping at the level of entire responses rather than at each token, thereby filtering out significantly “off-policy” samples during gradient estimation. This approach ensures better alignment with both sequence-level reward mechanisms and optimization objectives. Additionally, we introduce length normalization in $s_{i}(\theta)$ to minimize variance and to keep $s_{i}(\theta)$ values within a consistent numerical interval. Without this normalization, shifts in the probabilities of just a few tokens could cause substantial swings in the sequence-level importance ratio, which would necessitate different clipping thresholds depending on the response length. It is also worth noting that, because GSPO and prior methods (such as GRPO) define importance ratios differently, the clipping intervals they use generally differ by an order of magnitude.

\subsection{Gradient Analysis}
\label{subsec:gradient}

We can derive the gradient of the GSPO objective as follows (clipping is omitted for brevity):

\begin{align}
\nabla_{\theta} \mathcal{J}_\text{GSPO} (\theta)
=&\ 
\nabla_{\theta} \mathbb{E}_{ x \sim \mathcal{D},\, \{y_i\}_{i=1}^G \sim \pi_{\theta_\text{old}}( \cdot | x) }
\left[ 
\frac{1}{G} \sum_{i=1}^{G} s_{i}(\theta) \widehat{A}_{i}
\right] \\
= &\ 
\mathbb{E}_{ x \sim \mathcal{D},\, \{y_i\}_{i=1}^G \sim \pi_{\theta_\text{old}}( \cdot | x) }
\left[ 
\frac{1}{G} \sum_{i=1}^{G}
s_{i}(\theta) \widehat{A}_{i}
\cdot \nabla_{\theta} \log s_{i}(\theta)
\right] \\
=&\ 
\mathbb{E}_{ x \sim \mathcal{D},\, \{y_i\}_{i=1}^G \sim \pi_{\theta_\text{old}}( \cdot | x) }
\left[ 
\frac{1}{G} \sum_{i=1}^{G} 
\left( \frac{ \pi_{\theta} (y_i | x) }{ \pi_{\theta_\text{old}} (y_i | x)} \right)^{\frac{1}{|y_i|}} \widehat{A}_{i} 
\cdot \frac{1}{|y_i|} \sum_{t=1}^{|y_i|} \nabla_{\theta} \log \pi_{\theta} (y_{i,t} | x, y_{i,<t}) 
\right].
\label{equ:gspo_grad}
\end{align}

For comparison, the gradient of the GRPO objective is as follows (note that $\widehat{A}_{i,t} = \widehat{A}_{i}$):

\begin{align}
\nabla_{\theta} \mathcal{J}_\text{GRPO} (\theta) 
=&\ 
\nabla_{\theta} \mathbb{E}_{ x \sim \mathcal{D},\, \{y_i\}_{i=1}^G \sim \pi_{\theta_\text{old}}( \cdot | x) }
\left[ \frac{1}{G} \sum_{i=1}^{G} \frac{1}{|y_i|} \sum_{t=1}^{|y_i|} 
w_{i,t}(\theta) \widehat{A}_{i,t}
\right] \\
=&\
\mathbb{E}_{ x \sim \mathcal{D},\, \{y_i\}_{i=1}^G \sim \pi_{\theta_\text{old}}( \cdot | x) }
\left[ \frac{1}{G} \sum_{i=1}^{G} \widehat{A}_{i} 
\cdot \frac{1}{|y_i|} \sum_{t=1}^{|y_i|} 
\frac{ \pi_{\theta} (y_{i,t} | x, y_{i,<t}) }{ \pi_{\theta_\text{old}} (y_{i,t} | x,y_{i,<t})} 
\nabla_{\theta} \log \pi_{\theta} (y_{i,t} | x, y_{i,<t})  
\right].
\end{align}

The principal difference between GSPO and GRPO centers on \textit{the method used to weight the gradients of token log likelihoods}. In the case of GRPO, each token is assigned a weight based on its ``importance weight'' $\frac{ \pi_{\theta} (y_{i,t} | x, y_{i,<t}) }{ \pi_{\theta_\text{old}} (y_{i,t} | x,y_{i,<t})}$, resulting in potentially varied contribution per token. These weights are not uniformly distributed but instead can range within $(0, 1+\varepsilon]$ (when $\widehat{A}_{i} >0$) or $[1-\varepsilon, +\infty)$ (for $\widehat{A}_{i} < 0$), rendering them significant and non-negligible. Such disparity may accumulate over training iterations, ultimately inducing instability or unpredictable training dynamics. On the other hand, GSPO imposes uniform weighting across all tokens within a given response, thereby removing the instability that may arise from the non-uniform weighting in GRPO.

\subsection{GSPO-token: A Token-level Objective Variant}

In scenarios like multi-turn RL, we may desire a finer-grained advantage adjustment than the sequence level. To this end, we introduce a token-level objective variant of GSPO, namely \textbf{GSPO-token}, to allow token-wise advantage customization:

\begin{align}
\mathcal{J}_\text{GSPO-token}(\theta)
=
\mathbb{E}_{ x \sim \mathcal{D},\, \{y_i\}_{i=1}^G \sim \pi_{\theta_\text{old}}( \cdot | x) }
\left[ 
\frac{1}{G} \sum_{i=1}^{G} \frac{1}{|y_i|} \sum_{t=1}^{|y_i|} 
\min \left( s_{i,t}(\theta) \widehat{A}_{i,t},  \, \mathrm{clip} \left( s_{i,t}(\theta), 1 - {\varepsilon}, 1 + {\varepsilon} \right) \widehat{A}_{i,t} \right) 
\right],
\label{equ:gspo-token}
\end{align}

where

\begin{align}
s_{i,t}(\theta) 
= \mathrm{sg} \left[ s_{i}(\theta) \right]  \cdot \frac{ \pi_{\theta} (y_{i,t} | x, y_{i,<t}) }{ \mathrm{sg} \left[ \pi_{\theta} (y_{i,t} | x, y_{i,<t}) \right] },
\end{align}

and $\mathrm{sg}[\cdot]$ denotes only taking the numerical value but stopping the gradient, corresponding to the \texttt{detach} operation in PyTorch. The gradient of GSPO-token can be derived as:

\begin{align}
\nabla_{\theta} \mathcal{J}_\text{GSPO-token}(\theta)
=&\ 
\nabla_{\theta} \mathbb{E}_{ x \sim \mathcal{D},\, \{y_i\}_{i=1}^G \sim \pi_{\theta_\text{old}}( \cdot | x) }
\left[ 
\frac{1}{G} \sum_{i=1}^{G}
\frac{1}{|y_i|} \sum_{t=1}^{|y_i|} 
s_{i,t}(\theta) \widehat{A}_{i,t}
\right] \\
=&\ 
\mathbb{E}_{ x \sim \mathcal{D},\, \{y_i\}_{i=1}^G \sim \pi_{\theta_\text{old}}( \cdot | x) }
\left[ 
\frac{1}{G} \sum_{i=1}^{G} s_i (\theta) \cdot \frac{1}{|y_i|} \sum_{t=1}^{|y_i|} 
\widehat{A}_{i,t} \frac{ \nabla_{\theta} \pi_{\theta} (y_{i,t} | x, y_{i,<t}) }{ \pi_{\theta} (y_{i,t} | x, y_{i,<t}) }
\right] \\
=&\ 
\mathbb{E}_{ x \sim \mathcal{D},\, \{y_i\}_{i=1}^G \sim \pi_{\theta_\text{old}}( \cdot | x) }
\left[ 
\frac{1}{G} \sum_{i=1}^{G}  \left( \frac{ \pi_{\theta} (y_i | x) }{ \pi_{\theta_\text{old}} (y_i | x)} \right)^{\frac{1}{|y_i|}}
\cdot \frac{1}{|y_i|} \sum_{t=1}^{|y_i|}
\widehat{A}_{i,t} \nabla_{\theta} \log \pi_{\theta} (y_{i,t} | x, y_{i,<t})  
\right].
\label{equ:gspo-token_grad}
\end{align}

Observe that the expression $\frac{ \pi_{\theta} (y_{i,t} | x, y_{i,<t}) }{ \mathrm{sg} \left[ \pi_{\theta} (y_{i,t} | x, y_{i,<t}) \right] }$ evaluates to 1, which means that $s_{i,t}(\theta)$ is numerically identical to $s_{i}(\theta)$. A comparison between Equation~\eqref{equ:gspo} and~\eqref{equ:gspo-token}, as well as between Equation~\eqref{equ:gspo_grad} and~\eqref{equ:gspo-token_grad}, reveals that GSPO-token and GSPO produce equivalent values for the optimization objective, the clipping condition, and the theoretical gradient provided all tokens in the output $y_i$ share an identical advantage, namely $\widehat{A}_{i,t} = \widehat{A}_{i}$. However, GSPO-token further permits varying advantages across individual tokens, thereby offering increased flexibility.

\section{Experiments and Discussion}

\subsection{Empirical Results}

We conduct experiments utilizing a cold-start model that has been fine-tuned from Qwen3-30B-A3B-Base, and present both the training reward trajectories and the evaluation metrics for the model on the AIME'24 benchmark (average Pass@1 over 32 samples), LiveCodeBench (202410-202502, average Pass@1 over 8 samples), and CodeForces (measured via Elo Rating). Throughout reinforcement learning, rollout batches are subdivided into four mini-batches for performing gradient updates. For GSPO, the clipping bounds in Equation~\eqref{equ:gspo} are configured to 3e-4 (lower) and 4e-4 (upper). As a baseline, we use GRPO, assigning its left and right clipping intervals in Equation~\eqref{equ:grpo} to 0.2 and 0.27, respectively; these settings have been optimized to guarantee equitable comparisons. It is important to highlight that GRPO requires the Routing Replay training approach to reliably train MoE RL models, which we elaborate on in \S~\ref{subsec:routing_replay}, whereas \textit{GSPO eliminates the dependency on this approach}.

Figure~\ref{fig:results} demonstrates that training utilizing GSPO remains stable over the entire process. Notably, \textbf{GSPO consistently achieves higher performance by increasing the amount of training compute, frequently refreshing the query set, and extending the length of generated outputs}. Additionally, compared to GRPO, GSPO exhibits greater training efficiency, securing higher accuracy and benchmark results while using equivalent training compute and query resources. Lastly, GSPO has been effectively employed in RL training for state-of-the-art Qwen3 models, providing robust evidence of GSPO’s ability to harness the scaling benefits of RL for large language models.

\subsection{Curious Observation on Clipping Fractions}

One notable difference between GSPO and GRPO lies in their approach to clipping: GSPO applies clipping at the response level, while GRPO operates at the level of individual tokens. As depicted in Figure~\ref{fig:clipping}, this results in GSPO clipping a substantially greater proportion of tokens—by approximately two orders of magnitude—compared to GRPO, a pattern that persists even when the clipping ranges are modified. Despite the fact that GSPO excludes far more tokens from contributing to training (or gradient calculations), it nonetheless surpasses GRPO in terms of training efficiency. This seemingly paradoxical outcome—where eliminating a much larger subset of tokens results in improved training efficiency—suggests that gradient estimates at the token level in GRPO are prone to significant noise and are suboptimal for utilizing sampled data. Conversely, GSPO's reliance on full-sequence clipping enables it to provide a more robust and informative gradient signal, thereby enhancing learning effectiveness.

\subsection{Benefit of GSPO for MoE Training}
\label{subsec:routing_replay}

\paragraph{Background}
While reinforcement learning (RL) on dense models typically proceeds without major instability, MoE models pose distinctive challenges due to their sparse expert activation patterns. Notably, during our experiments utilizing the GRPO algorithm, we observed that MoE models experience significant \textit{expert-activation volatility}, which can impede successful RL convergence. Specifically, the set of experts chosen to generate identical responses may shift markedly after even a few gradient steps. Taking the 48-layer Qwen3-30B-A3B-Base model as an example, after each RL gradient update, approximately 10\% of the experts selected for the same rollout differ between the current policy $\pi_{\theta}$ and the previous policy $\pi_{\theta_\text{old}}$. This effect is exacerbated as model depth increases, resulting in severe fluctuations of the token-level importance ratios $w_{i,t}(\theta) = \frac{ \pi_{\theta} (y_{i,t} | x, y_{i,<t}) }{ \pi_{\theta_\text{old}} (y_{i,t} | x,y_{i,<t})}$, thereby undermining their reliability. As discussed in \S~\ref{sec:motivation} and~\ref{subsec:gradient}, these instabilities ultimately interfere with the standard convergence expected in RL training.

\paragraph{Our Previous Approach} In our earlier work, we addressed this issue using the \textbf{Routing Replay} training methodology. This method involves recording the set of experts activated under $\pi_{\theta_\text{old}}$ and then utilizing these identical routing patterns in $\pi_{\theta}$ during the calculation of the importance ratios $w_{i,t}(\theta) = \frac{ \pi_{\theta} (y_{i,t} | x, y_{i,<t}) }{ \pi_{\theta_\text{old}} (y_{i,t} | x,y_{i,<t})}$. By enforcing this, both $\pi_{\theta} (y_{i,t} | x, y_{i,<t})$ and $\pi_{\theta_\text{old}} (y_{i,t} | x,y_{i,<t})$ for each token $y_{i,t}$ operate on an identical set of activated network components, which stabilizes the computation of importance weights at the token level and guarantees consistent optimization within the activated subnetworks during gradient-based training steps. As illustrated in Figure~\ref{fig:routing_replay}, Routing Replay proves to be a critical mechanism for achieving smooth and reliable convergence in GRPO training procedures involving MoE architectures.

\paragraph{Benefit of GSPO} While Routing Replay supports effective convergence during GRPO training for MoE models, its strategy of reapplying previous routing decisions results in higher memory usage, increased communication costs, and can restrict the practical capacity of the MoE system. In comparison, as depicted in Figure~\ref{fig:results}, GSPO removes reliance on Routing Replay altogether and can perform traditional computation of importance ratios $s_i(\theta)$, achieving stable convergence and optimization. The central idea behind GSPO is its exclusive focus on the overall sequence likelihood (i.e., $\pi_{\theta} (y_i | x)$), exhibiting insensitivity to the likelihoods of specific tokens (i.e., $\pi_{\theta} (y_{i,t} | x, y_{i,<t})$). Since the MoE architecture retains robust language modeling abilities throughout, the sequence likelihood remains largely stable. Ultimately, GSPO addresses the instability in expert activation found in MoE models at its root, making workarounds like Routing Replay unnecessary. As a result, this method not only streamlines and fortifies the training procedure but also ensures the model can fully utilize its inherent capacity without imposed limits.

\subsection{Benefit of GSPO for RL Infrastructure}

Due to variations in numerical precision between training systems such as Megatron and inference systems like SGLang and vLLM, it is common practice to recalculate the likelihoods of generated responses under the old policy $\pi_{\theta_\text{old}}$ using the training engine. Notably, GSPO focuses on optimizing based on sequence-level likelihoods as opposed to token-level ones; this broader granularity tends to be significantly less sensitive to discrepancies in precision. As a result, GSPO enables direct use of the inference engine’s likelihood outputs during optimization, eliminating the necessity of recalculating them with the training engine. This approach proves particularly advantageous for use cases such as partial rollout, multi-turn reinforcement learning, and frameworks where training and inference are handled separately.

\section{Conclusion}

We introduce Group Sequence Policy Optimization (GSPO), a novel reinforcement learning approach tailored for training large language models. Rooted in the concept of importance sampling, GSPO constructs importance ratios using sequence likelihood, applying sequence-level procedures for clipping, rewarding, and policy optimization. In comparison to GRPO, GSPO provides significant advantages in terms of training stability, efficiency, and overall performance. The algorithm proves particularly effective in large-scale RL scenarios involving MoE models, underpinning the substantial advancements observed in the latest Qwen3 models. By establishing GSPO as a robust, scalable foundation, we aim to further expand RL capabilities and anticipate transformative progress in artificial intelligence.

\end{document}